% ---------------------------------------------------------------- 3 problem, 4 method
\section{From pixels to a knot diagram}\label{sec:method}

\paragraph{Diagram model.} A diagram is a set of closed strands (polylines in the
image plane) and crossings. Each crossing has two passages: an over-strand and an
under-strand. Its sign is the right-hand rule $\operatorname{sign}(\mathbf{t}_o\times\mathbf{t}_u)$
on the strand tangents, and the diagram's PD code follows by walking each strand.
Identification is then exact: a hyperbolic exterior is matched by an isometry
certificate, and knot Floer homology settles the remaining cases, including the
unknot and chirality. An image fixes the diagram only if every crossing is visible
and its \ou{} can be seen. A silhouette shows crossings but not \ou{}, so it fixes
the projection but not the knot.

\paragraph{Input.} The image is resized to $256\times256$ and given four channels:
RGB and an FFT-cleaned grey channel (mirror-padded homomorphic filtering, which
removes slow illumination, followed by Wiener shrinkage).

\paragraph{Networks.} Any network that outputs the maps below can be used. We study two.
(a) A NanoTransformer (0.86M parameters): a convolutional stem (stride 4) feeds a $16\times16$ token grid plus
one count token into four transformer blocks (width 128, 4 heads). Global attention
lets evidence from both sides of a crossing meet. A light U-Net decoder with
stride-4 and stride-2 skips returns to a $128\times128$ grid, where five heads predict
(\cref{fig:method}):
\begin{itemize}
\item a soft cable mask $m$ (fraction of each cell covered);
\item a crossing heatmap $h$, trained with a centre-point focal loss;
\item the over-strand axis $\mathbf{o}=(\cos2\alpha,\sin2\alpha)$ near each crossing;
\item the strand axis elsewhere, used to pair arms;
\item a centre-line ridge $c=\exp(-d^2/2\sigma^2)$, where $d$ is the distance to the
  nearest strand centre-line.
\end{itemize}
A full-resolution branch (input-level convolutions fused with the decoder) predicts
the same ridge on a $256\times256$ grid. On the coarse grid, two centre-lines one
\cw{} apart are only 1.5--3 cells apart and merge. Axial (doubled-angle) targets make
a strand and its reverse identical, which is all an image can show. Dihedral
augmentations transform the axial targets consistently. (b) A pair of standard U-Nets
(1.6M parameters each, 5 levels from $256^2$ to $16^2$): one predicts the mask and the
centre-lines, the other the crossing heatmap with the over-strand axis as a per-keypoint
attribute, as in CenterNet~\cite{zhou2019objects}. (c) A \emph{UNetFormer} (2.9M parameters), a single
network that predicts every map. It has a U-Net encoder and decoder with skip connections at all five
resolutions, as in (b), and four transformer blocks (width 192, 6 heads) with a count token on the
$16\times16$ grid at the bottleneck, as in TransUNet~\cite{chen2021transunet}. The transformer's output
is added back to the bottleneck through a zero-initialised projection, so training starts from a
plain U-Net.

\paragraph{Parser.} (i) Crossings are heatmap peaks above 0.35, refined to sub-cell
position. (ii) The skeleton is taken from the thresholded ridge. Specks below
$(0.6w)^2$ pixels are removed, where $w$ is the estimated \cw. Short gaps between
facing ends are bridged, and the skeleton is cut on a disc around each crossing,
leaving arcs whose ends are the crossing's arms. On cluttered backgrounds the networks draw short
centre-line pieces on background texture. Of 134 free ends that made the NanoTransformer abstain on
the benchmark's unseen backgrounds, 128 lay off the cable. A separate piece that carries no
detected crossing and is open (its skeleton has an end) cannot belong to a closed diagram, so it is
removed. Closed loops without crossings, such as an unknot or an unlinked component, are kept.
(iii) Of the three ways to split
four arms into two strands, we take the one whose strands run straightest. (iv) The
strand whose axis agrees best with $\mathbf{o}$ is on top. Its probability is
$p=\tfrac12+\big(\sigma(\kappa\,\Delta)-\tfrac12\big)\,\rho^{\gamma}$, where $\Delta$ is the
margin between the two strands, $\kappa$ is fitted on the training split, and $\rho$
is the mean resultant length of $\mathbf{o}$ across eight test-time views
($\gamma=2$, chosen on a development set). Where the views disagree, as on a
silhouette, $\rho\to0$ and $p\to\tfrac12$. (v) Tracing arcs through crossings yields
oriented strands and passages, i.e.\ a diagram.

\paragraph{Gap evidence (optional).} View agreement cannot detect \ou{} evidence that every
view misreads in the same way. The picture can be checked directly: the strand on top keeps
its outline through the crossing, and the strand underneath loses it. For each strand through a
crossing we sample, every half pixel of arc length $t$ in the 512-pixel image, the contrast
across both of its edges at $\pm w/2$ on the normal. The edge is \emph{missing} where this
contrast is below $\tau=\GapTau$ times its median on the free part of the strand. The strand's
interruption is a Lebesgue measure,
\[ G=\tfrac1w\,\big|\{t\in[-T,T]:\text{edge missing at }t\}\big|,\qquad T=\tfrac{w}{2\sin\theta}+w,\]
where $\theta$ is the crossing angle. $D=G_{\text{under}}-G_{\text{over}}$ is the picture's
evidence for the network's reading. It enters the log-odds,
$\operatorname{logit}p'=\operatorname{logit}p+\lambda D$, and the reading flips if
$p'<\tfrac12$. $\lambda=\GapLambda$ maximises the likelihood of \GapDevCross{} matched
crossings on development images; no test image is involved.

\paragraph{Minkowski attention and a Lebesgue loss.} We test three changes to the
NanoTransformer, each added to the previous one.
(1)~\emph{Lebesgue (layer-cake) loss.} The parser thresholds the mask at levels
$t\in[0.15,0.85]$ and the ridges at 0.5. By the layer-cake formula, the mean over
$t\sim U[a,b]$ of the area of $\{\hat p>t\}\triangle\{m>t\}$ equals
$\mathbb E\,|\mathrm{clip}(\hat p,a,b)-\mathrm{clip}(m,a,b)|/(b-a)$. We add this term with weight 1
for the mask ($[0.15,0.85]$) and both ridges ($[0.3,0.7]$). It penalises exactly the pixels
whose thresholded value depends on the level.
(2)~\emph{Minkowski position bias.} Each head adds $-\alpha_h\|\mathbf r_i-\mathbf r_j\|_{p_h}$ to
its logits, where $\mathbf r$ are token positions on the $16\times16$ grid. Both
$p_h=1+\operatorname{softplus}(\rho_h)$ (initially 2) and $\alpha_h$ (initially 0.2 down to 0.01
across heads) are learned per block. The $\ell_p$ distance is invariant under the eight flips and
quarter turns of the square, the group used for test-time views.
(3)~\emph{Minkowski content attention.} The logits are
$-t_h\|\mathbf q_i-\mathbf k_j\|_p^p/(p\sqrt d)$, with $p=1$ on two heads and $p=2$ on the other
two. For $p=2$ this is dot-product attention less a key-norm term. A learned $p$ would require
a per-coordinate $|\mathbf q-\mathbf k|^p$ tensor, about 400 times slower on a CPU.

\section{Consistency checks and abstention}\label{sec:gates}
The system reports a type only if all of the following hold; otherwise it abstains
and names the failed check.
\begin{enumerate}
\item \textbf{Euler count.} The projection of $k$ closed curves with $n$ crossings is a
  4-valent planar graph, so the ridge band has $n+\pi$ bounded faces, where $\pi$ is
  the number of connected pieces. A missed crossing adds an unexplained face.
\item \textbf{Four arms.} Every detected crossing has exactly four arms. No skeleton
  branch point lies away from a detected crossing.
\item \textbf{Closure.} The traversal closes. Free ends require a recorded closure
  (a fixture, or ends joined), as for open cables.
\item \textbf{Two resolutions agree.} The diagram is parsed from both ridges. If both
  pass, they must give the same type.
\item \textbf{Type probability.} Uncertain crossings are enumerated over their two
  readings, weighted by $p$, and the most likely type must reach probability 0.95.
\end{enumerate}
